% ------------------------------------------------------------------ cover (page 1)
\thispagestyle{blank}
\begin{tikzpicture}[remember picture,overlay]
  \fill[sun!35] ($(current page.south west)+(0.4in,0.4in)$) rectangle ($(current page.north east)-(0.4in,0.4in)$);
  % rays
  \begin{scope}
  \clip ($(current page.south west)+(0.4in,0.4in)$) rectangle ($(current page.north east)-(0.4in,0.4in)$);
  \foreach \a in {0,15,...,345}{
    \fill[sun!70,opacity=0.5] ($(current page.center)+(0,0.8in)$) -- ++(\a:9in) -- ++(\a+7.5:9in) -- cycle;
  }
  \end{scope}
  \fill[sun!35,opacity=0.9] ($(current page.south west)+(0.4in,0.4in)$) rectangle ($(current page.south east)+(-0.4in,2.4in)$);
  \node[star,star points=18,star point ratio=0.88,fill=tangerine,draw=ink,line width=1.6pt,minimum size=2.7cm,
        font=\sffamily\bfseries,text=white,align=center,inner sep=0]
     at ($(current page.north west)+(1.55in,-1.7in)$) {\Large SCIENCE\\[-1pt]\large Book 8};
  \node[anchor=north,align=center,font=\sffamily\bfseries,text=ink,text width=5.7in]
     at ($(current page.north)+(0.4in,-1.45in)$) {\fontsize{32}{42}\selectfont Where\\[8pt] the Energy\\[8pt] Goes};
  % ramp, ball, bounce and warmth
  \draw[line width=2pt,color=ink] ($(current page.center)+(-2.4in,0.1in)$) -- ($(current page.center)+(-0.3in,-1.2in)$) -- ($(current page.center)+(2.4in,-1.2in)$);
  \fill[tangerine,draw=ink,line width=1.2pt] ($(current page.center)+(-2.1in,0.33in)$) circle (0.16in);
  \draw[-{Stealth},line width=2.4pt,color=sky!80!black] ($(current page.center)+(-1.7in,0.05in)$) -- ($(current page.center)+(-1.0in,-0.35in)$);
  \draw[-{Stealth},line width=2.4pt,color=leaf!70!black] ($(current page.center)+(0.0in,-1.1in)$) -- ($(current page.center)+(1.2in,-1.1in)$);
  \foreach \x in {1.5,1.8,2.1}{
    \draw[-{Stealth},line width=2.2pt,color=tomato,decorate,decoration={snake,amplitude=1.5pt,segment length=7pt}] ($(current page.center)+(\x in,-1.1in)$) -- ($(current page.center)+(\x in,-0.4in)$);
  }
  \node at ($(current page.center)+(1.6in,-2.35in)$) {\scalebox{3.1}{\multixbot[wow]}};
  \node[anchor=center,fill=white,draw=ink,line width=1.2pt,rounded corners=6pt,font=\sffamily\bfseries\small,
        text width=2.2in,align=center,inner sep=6pt]
    at ($(current page.center)+(-1.15in,-2.55in)$) {Lift it. Roll it. Warm it. Keep the total.\\[2pt]
    \footnotesize Puzzles, projects and Multix the guide robot};
  \node[font=\sffamily\bfseries\large,text=ink] at ($(current page.south)+(0,0.8in)$) {Ages 9 to 11};
\end{tikzpicture}
\clearpage

% ------------------------------------------------------------------ inside front (page 2)
\thispagestyle{plain}
\begin{center}
{\sans\bfseries\large This book belongs to}\quad\wline[6cm]
\end{center}
\vspace{0.1cm}

{\sans\bfseries\LARGE\color{ink} Hello. I am Multix.}\par
\begin{multix}[wow]\multixbot[wow]\tcblower
I am a guide robot. I read every instruction very carefully, which is why I sometimes get things wrong.
When I make a mistake, see whether you can catch it. I never give away a puzzle answer.
\end{multix}

\subhead{How this book works}
Every book in this series asks the same few questions, so once you have read this page you will know your way around.

\begin{itemize}[leftmargin=1.6em,itemsep=0.45em]
\item[\textcolor{leaf}{\ding{108}}] \textbf{Try it} boxes tell you what to do with a ruler, a cup, some string or a clock. You need nothing from a shop. Everything is at home or at school.
\item[\textcolor{sun!80!black}{\ding{108}}] \textbf{What did we keep?} After you simplify something, such as keeping the same conditions in a test or averaging several throws, you say what stayed and what was set aside. Simplifying is fine. Forgetting what you set aside is not.
\item[\textcolor{tangerine}{\ding{108}}] \textbf{What makes this correction trustworthy?} When you fix or improve an answer, a box says what you checked, what the check gives you reason to trust, and what it does \emph{not} show.
\item[\textcolor{sky}{\ding{108}}] \textbf{Circles} show how a number was made. Page~\pageref{circlecard} explains how to read them, and Book~1 introduced them.
\item[\textcolor{grape}{\ding{108}}] \textbf{Puzzles} come at the end of each chapter. Try them before you look at the answers on page~\pageref{answers}.
\item[\textcolor{tomato}{\ding{108}}] \textbf{Slop check} is a page near the end. Slop is hasty, gaudy work that forgot to take care. You will learn to spot it and to fix it.
\end{itemize}

\begin{multix}[hmm]\multixbot[hmm]\tcblower
Here is a fact about me. I watched a ball bounce lower and lower until it stopped. I announced that its energy had been destroyed. Chapter~1 will explain why I should be embarrassed.
\end{multix}
\clearpage

% ------------------------------------------------------------------ contents and kit (page 3)
\thispagestyle{plain}
{\sans\bfseries\LARGE\color{ink} In this book}\par\vspace{0.2cm}
{\sans\large
\begin{tabular}{@{}p{0.9cm}p{7.6cm}r@{}}
\textcolor{tangerine}{\bfseries 1} & Energy You Can Lift & \pageref{ch1}\\[3pt]
\textcolor{tangerine}{\bfseries 2} & Where a Rolling Ball Goes & \pageref{ch2}\\[3pt]
\textcolor{tangerine}{\bfseries 3} & Warm Things Cool Down & \pageref{ch3}\\[3pt]
\textcolor{tangerine}{\bfseries 4} & Machines Trade Distance & \pageref{ch4}\\[3pt]
\textcolor{tangerine}{\bfseries 5} & Keeping a Running Total & \pageref{ch5}\\[3pt]
\textcolor{tangerine}{\bfseries 6} & Finding Missing Energy & \pageref{ch6}\\[3pt]
\textcolor{sky}{\bfseries $\star$} & Slop check & \pageref{slop}\\[3pt]
\textcolor{sky}{\bfseries $\star$} & Answers & \pageref{answers}\\[3pt]
\textcolor{sky}{\bfseries $\star$} & Words, and notes for grown-ups & \pageref{notes}\\
\end{tabular}}

\begin{tryit}{Gather your kit}
You will use these things across the six chapters. Everything is found at home or at school. Measuring comes from Book~1, fair testing from Book~2, forces from Book~4 and the ledger from Books~5 and~6, and all of them come back here.
\begin{multicols}{2}
\begin{itemize}[leftmargin=1.2em]
\item a bouncy ball
\item a board and some books
\item a ruler and a tape measure
\item two cups, one with a lid
\item a kitchen thermometer in \textdegree C
\item a clock or timer
\item a rubber band and a bag of coins
\item dominoes or small blocks
\item a pencil and notebook
\item a hard floor and a rug
\end{itemize}
\end{multicols}
\end{tryit}
\adult{Use warm tap water only, never boiling water from a kettle or stove. Drop things only from low heights and keep ramps low and away from stairs. Keep small balls and coins away from little children and pets. Nothing here needs a flame or mains electricity.}

{\sans\bfseries\large\color{sky!70!black} Units in this book}\par
We use the metric system: centimetres (cm), metres (m), seconds (s) and degrees Celsius (\textdegree C). Energy is hard to measure at home, so we keep a ledger in made-up \textbf{units}: for the same ball, 1~unit is 1~cm of lift.
There are 10~mm in 1~cm and 100~cm in 1~m. If you are used to inches, one inch is about 2.5~cm. We will use that only in side notes.
\clearpage

% ------------------------------------------------------------------ circle card (page 4)
\thispagestyle{plain}\label{circlecard}
{\sans\bfseries\LARGE\color{ink} Reading the circles}\par\vspace{0.1cm}
A circle holds a group of things together. There are three things you can do with circles in this series.

\begin{center}
\begin{tikzpicture}[scale=1.0]
  % POP
  \node[font=\sffamily\bfseries,text=tomato] at (0,2.6) {1. POP};
  \node[bigcir,minimum size=2.2cm] (p) at (0,0.8) {};
  \foreach \x/\y/\t in {-0.35/1.15/2.1,0.4/1.2/1.9,-0.45/0.5/2.0,0.4/0.45/2.2}{\node[cir,inner sep=1pt,font=\sffamily\bfseries\scriptsize] at (\x,\y) {\t};}
  \node[valcir] (pv) at (0,-1.5) {2.05};
  \draw[popA] (p) -- node[right,note] {pop} (pv);
  % MOVE
  \node[font=\sffamily\bfseries,text=sky!70!black] at (4.0,2.6) {2. MOVE};
  \node[cir] (m1) at (3.0,0.8) {12}; \node[cir] (m2) at (5.0,0.8) {7};
  \draw[moveA] (m1) -- node[above,note] {5 beans} (m2);
  \node[note,anchor=east] at (2.35,0.8) {before};
  \node[cir] at (3.0,-1.5) {7}; \node[cir] at (5.0,-1.5) {12};
  \node[note,anchor=east] at (2.35,-1.5) {after};
  % WRAP
  \node[font=\sffamily\bfseries,text=tangerine!80!black] at (8.0,2.6) {3. WRAP};
  \node[wrapcir,minimum size=2.5cm] (w) at (8.0,0.5) {};
  \node[valcir,font=\sffamily\bfseries\scriptsize] at (8.0,0.9) {14.5};
  \node[note] at (8.0,0.05) {$-\,2.0$};
  \node[valcir] at (8.0,-1.5) {12.5};
\end{tikzpicture}
\end{center}

\begin{description}[leftmargin=0pt,itemsep=0.4em,style=unboxed]
\item[\textcolor{tomato}{\sans\bfseries Pop}] Squash a group down to one value that sums it up, such as an average. A pop is a \emph{summary}, so it always loses something. Every time we pop, we say what the summary keeps and what it loses.
\item[\textcolor{sky!70!black}{\sans\bfseries Move}] An arrow that carries something from one circle to another, such as beans or water. Nothing is summarized. You will see more of moves in later books.
\item[\textcolor{tangerine!80!black}{\sans\bfseries Wrap}] When you correct something, draw a new dashed circle around the old one and show what changed. The old work stays visible, so anyone can check it.
\end{description}

\subhead{Energy moves between circles}
\begin{center}
\begin{tikzpicture}
  \node[bigcir,minimum size=2.2cm] (a) at (0,0) {}; \node[font=\sffamily\scriptsize,align=center] at (0,0) {height\\40};
  \node[cir] (b) at (4.8,1.0) {30}; \node[cir] (c) at (4.8,-1.0) {10};
  \draw[moveA] (a) -- node[above,note] {motion} (b);
  \draw[moveA] (a) -- node[below,note] {warmth} (c);
  \node[note,anchor=west,text width=3.8cm,align=left] at (5.6,0) {Energy moves from one circle to others. The total, $30+10=40$, stays the same.};
\end{tikzpicture}
\end{center}
\begin{multix}[ok]\multixbot[ok]\tcblower
A pop is not a lie, but it is a shortcut. Any time I tell you a pop, I will tell you what it threw away.
\end{multix}
